\begin{center}
  \section*{Homework 04 - 2.6, 2.7}
  Due Mar 6 \\
  Uzair Hamed Mohammed
\end{center}

% 6, 16, 18, 26, 38, 58
% 6, 12, 20

\subsection*{2.6 Combinatorics}
\begin{enumerate}
  \item[6.] A fast-food restaurant offers customers a choice of eight
    toppings that can be added to a hamburger. How many different
    hamburgers can be ordered?

    \underline{Sol}: \(\boxed{2^8 = 256}\) different hamburgers can be
    ordered since each topping can be added for two choices of each topping.

  \item[16.] Monica’s vacation plans require that she fly from
    Nashville to Chicago to Seattle to Anchorage. According to her
    travel agent, there are three available flights from Nashville to
    Chicago, five from Chicago to Seattle, and two from Seattle to
    Anchorage. Assume that the numbers of options she has for return
    flights are the same. How many round-trip itineraries can she schedule?

    \underline{Sol}: On her departing trip, she can schedule \(3
    \times 5 \times 2 = 30\) itineraries. On her way back, she can
    schedule the same number. So, in total, she can schedule \(\boxed{30^2}\)
    different round-trip itineraries.

  \item[18.] How many ways can a set of four tires be put on a car if
    all the tires are interchangeable? How many ways are possible if
    two of the four are snow tires?

    \underline{Sol}: \(_2 P_0 = 1\) if all the tires are
    interchangeable, and we can select the location of the first snow
    tire in four ways and the location of the second snow tire in
    three ways to make a total of 12 ways in the second part.

  \item[26.] A new horror movie, Friday the 13th, Part X, will star
    Jason’s great-grandson (also named Jason) as a psychotic trying
    to dispatch (as gruesomely as possible) eight camp counselors,
    four men and four women. (a) How many scenarios (i.e., victim
    orders) can the screenwriters devise, assuming they want Jason to
    do away with all the men before going after any of the women? (b)
    How many scripts are possible if the only restriction imposed on
    Jason is that he save Muffy for last?

    \begin{enumerate}
      \item \underline{Sol}: There are \(4!\) ways to order the four
        men and \(4!\)
        ways to order the four women, meaning there are \(\boxed{4!
        \times 4! = 576}\) scenarios where the men are dispatched
        before the women.

      \item \underline{Sol}: Since Muffy is fixed in the last
        position, there are
        \(7!\) ways to arrange the remaining seven camp counselors,
        resulting in \(\boxed{7! = 5040}$ possible scripts
      \end{enumerate}

    \item[38.] How many ways can the letters in the word
      \[
        S L U M G U L L I O N
      \]
      be arranged so that the three L’s precede all the other consonants?

      \underline{Sol}: There are \(\binom{11}{4} = 330\) ways to
      choose the positions for the vowels and \(\frac{4!}{2!} = 12\)
      ways to arrange the set \(\{U, U, I, O\}\). Since the three
      \(L\)s must precede the other four consonants, they must occupy
      the first three consonant positions in the sequence, leaving
      \(4! = 24\) ways to arrange the remaining consonants (\(S, M,
      G, N\)), resulting in a total of \(\boxed{330 \times 12 \times
      24 = 95,040}\) arrangements.

    \item[58.] Prove that
      \[
        \binom{n+1}{k} = \binom{n}{k} + \binom{n}{k-1}
      \]

      \underline{Sol}: Expanding the right-hand side with factorials
      gives $\frac{n!}{k!(n-k)!} + \frac{n!}{(k-1)!(n-k+1)!}$. By
      using a common denominator of $k!(n-k+1)!$, we sum the
      numerators to get $n!(n-k+1+k) = n!(n+1)$, which simplifies to
      $\frac{(n+1)!}{k!(n+1-k)!} = \boxed{\binom{n+1}{k}}$.
  \end{enumerate}

  \subsection*{2.7 Combinatorial Probability}

  \begin{enumerate}
    \item[6.] A committee of fifty politicians is to be chosen from
      among one hundered US senators. If the selection is done at
      random, what is the probability that each state will be represented?

      \underline{Sol}: We can choose 50 politicians in
      \(\binom{100}{50}\) ways. For each state, we can choose one
      representative in \(\binom{2}{1} = 2\) ways. In the end, we get
      probability:

      \[
        \boxed{P = \frac{2^{50}}{\binom{100}{50}}}
      \]

    \item[12.] If the letters in the phrase
      \[
        A \ R O L L I N G \ S T O N E \ G A T H E R S \ N O \ M O S S
      \]
      are arranged at random, what are the chances that not all the Ss
      will be adjacent?

      First, we have the following amount of unique permutations of
      all of the 26 letters:

      \[
        \frac{26!}{4!4!3!2!2!2!2!2!2!}
      \]

      Among these permutations, one is a subset contains all of the
      Ss adjacent. We can order all of these in
      \[
        \boxed{\frac{(26 - 4 + 1)!}{1!4!3!2!2!2!2!2!2!}}
      \]

      different ways.

    \item[20.] Five cards are dealt from a standard 52-card deck. What
      is the probability that the sum of the faces on the five cards is
      48 or more?

      \underline{Sol}: In a standard 52-card deck, if we assign the
      ranks Ace through King the values (1, 2, \(\dots\), 13), the total
      number of ways to be dealt five cards is \(\binom{52}{5} =
      2,598,960\). To find the number of hands where the sum of the
      faces is 48 or more, we sum the combinations of ranks that
      satisfy this condition (considering there are 4 of each rank).
      There are (158,472) such combinations, making the probability:
      $$P(\text{Sum} \geq 48) = \frac{158,472}{2,598,960} \approx
      \boxed{0.061}$$
  \end{enumerate}
